\chapter{Methodology}
\label{ch:methodology}

This chapter describes the experimental setup. We first formalize the MEXA pipeline as implemented in our codebase (Section~\ref{sec:mexa-pipeline}), then describe the evaluated models (Section~\ref{sec:models}), the parallel corpora and their subsets (Section~\ref{sec:datasets}), and the experimental settings, analysis methods, and compute environment (Section~\ref{sec:experimental-settings}). The methodology follows \citet{kargaran2025mexa} for the reproduction experiments and extends it along four axes: new decoder families, mixture-of-experts architectures, encoder and embedding models, and non-English pivot languages.

% ------------------------------------------------------------
\section{The MEXA Pipeline}
\label{sec:mexa-pipeline}

MEXA (Multilingual Evaluation via Cross-Lingual Alignment) estimates the multilingual coverage of a language model~$m$ from parallel sentences alone. We first describe the baseline four-stage evaluation pipeline as introduced by \citet{kargaran2025mexa} (Section~\ref{subsec:original-pipeline}), and then detail our proposed methodology extensions (Section~\ref{subsec:pipeline-extensions}).

% ..............................................
\subsection{Standard MEXA Pipeline}
\label{subsec:original-pipeline}

The reference MEXA framework operates in four sequential stages: sentence-embedding extraction, layer-wise similarity computation, alignment scoring, and aggregation, as schematically overviewed in Figure~\ref{fig:mexa-pipeline-diagram}.

\begin{figure}[H]
\centering
\includegraphics[width=0.98\textwidth]{figures/fig_mexa_pipeline_diagram.pdf}
\caption{\textbf{Schematic Overview of the MEXA Evaluation Pipeline}. The 4-stage pipeline extracts layer-wise hidden states, pools them with position weighting, computes the target-to-pivot similarity matrix $C^{(l)}$, applies bidirectional retrieval (P@1) filtering, and aggregates layer-wise scores into $\mu_{\text{Max}}$ and $\mu_{\text{Mean}}$.}
\label{fig:mexa-pipeline-diagram}
\end{figure}

\paragraph{Stage 1: Layer-wise sentence embeddings.}
Let $\mathcal{S} = \{s_1, \dots, s_n\}$ be a set of $n$ parallel sentences available in a target language $L_1$ and a pivot language $L_2$ (English by default, but configurable). Each sentence $s_i$ is passed through the model, and hidden states $h_{i,t}^{(l)}$ are collected for every token $t \in \{1, \dots, T\}$ at each layer $l \in \{0, \dots, L\}$ (where $l=0$ represents the input embedding layer). For a given layer $l$, the token hidden states $\{h_{i,t}^{(l)}\}_{t=1}^T$ are reduced to a single sentence embedding $e_i^{(l)}$ using one of two strategies. The first is a position-weighted average, where tokens receive weights linearly proportional to their position index $t$—so that later tokens in the sequence receive larger weights than earlier ones—as shown in Equation~\eqref{eq:weighted-avg}:
\begin{equation}
    e_i^{(l)} \;=\; \sum_{t=1}^{T} w_t\, h_{i,t}^{(l)},
    \qquad
    w_t \;=\; \frac{t}{\sum_{k=1}^{T} k},
    \label{eq:weighted-avg}
\end{equation}
where $T$ is the number of non-padding tokens. This is the primary setting in all reported results. The second is last-token pooling, which uses the hidden state of the final non-padding token at layer $l$, $e_i^{(l)} = h_{i,T}^{(l)}$. This position has attended to the entire sentence under causal attention, and is used as a robustness check. \citet{kargaran2025mexa} found that position-weighted average pooling consistently yields higher and more stable alignment scores than last-token pooling across causal decoders, as aggregating across all tokens mitigates sensitivity to single-token noise while preserving causal context accumulation.

For encoder-only models (Section~\ref{subsec:encoder-models}), whose attention is bidirectional, the same two reductions apply to the encoder hidden states, ensuring the pipeline is architecture-agnostic. Extracted embeddings are saved per language and per layer; consequently, alignment scoring (lightweight matrix calculations) is decoupled from embedding extraction (expensive GPU forward passes), enabling efficient re-evaluation of scoring variants without re-running model inference.

\paragraph{Stage 2: Similarity matrix.}
For each layer~$l \in \{0, \dots, L\}$ of language model~$m$, a pairwise cosine similarity matrix $C^{(l)} \equiv C(L_1, L_2, m, l) \in \mathbb{R}^{n \times n}$ is constructed across the full corpus of $n$ parallel sentences between target language $L_1$ and pivot language $L_2$, as shown in Equation~\eqref{eq:similarity-matrix}:
\begin{equation}
    C^{(l)} \;=\;
    \begin{pmatrix}
        c_{11}(l) & c_{12}(l) & \dots & c_{1n}(l) \\
        c_{21}(l) & c_{22}(l) & \dots & c_{2n}(l) \\
        \vdots    & \vdots    & \ddots & \vdots    \\
        c_{n1}(l) & c_{n2}(l) & \dots & c_{nn}(l)
    \end{pmatrix},
    \qquad
    c_{ij}(l) \;=\; \cos\!\left(e_{i, L_1}^{(l)},\, e_{j, L_2}^{(l)}\right),
    \label{eq:similarity-matrix}
\end{equation}
where entry $c_{ij}(l)$ denotes the cosine similarity between target sentence~$i \in L_1$ and pivot sentence~$j \in L_2$. The diagonal entries $c_{ii}(l)$ correspond to true parallel sentence pairs, while all off-diagonal entries $c_{ij}(l)$ ($i \neq j$) represent non-parallel distractors.

\paragraph{Stage 3: Alignment score.}
As discussed in Section~\ref{subsec:alignment-challenges}, absolute cosine values are unreliable in anisotropic LLM embedding spaces. MEXA therefore scores alignment with a \emph{bidirectional retrieval} (precision-at-1) criterion, as defined in Equation~\eqref{eq:mexa}: a parallel sentence pair~$(i, i)$ counts as aligned at layer~$l$ only if its true parallel similarity $c_{ii}(l)$ strictly exceeds every distractor entry in both its row ($c_{ij}(l)$, $j \neq i$) and its column ($c_{ji}(l)$, $j \neq i$),

\begin{equation}
    \mu C(L_1, L_2, m, l)
    \;=\;
    \frac{1}{n}
    \sum_{i=1}^{n}
    \mathbb{1}\!\left(
        c_{ii}(l) \;>\; \max_{j \neq i}\,\{\, c_{ij}(l),\; c_{ji}(l) \,\}
    \right).
    \label{eq:mexa}
\end{equation}

The score lies in $[0, 1]$. The expected value under random embeddings is approximately $1/n^2$ per pair (both the row and the column condition must hold by chance), so with $n = 100$ sentences a random baseline scores essentially zero. This validates the use of small parallel samples.

\paragraph{Stage 4: Aggregation across layers.}
The layer-wise scores are summarized into two model-level statistics, as shown in Equation~\eqref{eq:mexa-aggregation}:
\begin{equation}
    \mu_{\text{Max}}(L_1) = \max_{l} \; \mu C(L_1, L_2, m, l),
    \qquad
    \mu_{\text{Mean}}(L_1) = \frac{1}{L} \sum_{l=1}^{L} \mu C(L_1, L_2, m, l).
    \label{eq:mexa-aggregation}
\end{equation}
$\mu_{\text{Max}}$ captures the best alignment a model reaches at any depth, typically in the middle-to-late layers, consistent with the pivot hypothesis. $\mu_{\text{Mean}}$ rewards models that sustain alignment across many layers.

% ..............................................
\subsection{Proposed Pipeline Extensions}
\label{subsec:pipeline-extensions}

To extend the baseline MEXA framework, we introduce two methodological extensions designed to isolate representation geometry and test structural pivot assumptions.

\paragraph{Mean-centered alignment scoring.}
To separate semantic alignment from the language identity signal (Section~\ref{subsec:alignment-challenges}), we propose a centered variant of Equation~\eqref{eq:mexa}. This extension is inspired by \citet{schuster2019crosslingual}, who demonstrated that multilingual contextual representations contain a dominant language-specific centroid offset that can artificially separate language clusters at the expense of semantic alignment.

Before building the similarity matrix at layer~$l$, each language's embeddings are shifted by their sample mean vector, yielding the centered sentence embedding $\tilde{e}_i^{(l)}$ defined in Equation~\eqref{eq:mean-centering}:
\begin{equation}
    \tilde{e}_i^{(l)} \;=\; e_i^{(l)} - \frac{1}{n}\sum_{k=1}^{n} e_k^{(l)},
    \label{eq:mean-centering}
\end{equation}
applied independently to the target language $L_1$ and pivot language $L_2$. As illustrated geometrically in Figure~\ref{fig:meancentering-geometry} via a two-dimensional Principal Component Analysis (PCA) projection of the sentence embedding space, subtracting language-specific means centers both vector manifolds at the origin, eliminating the dominant language-identity vector $\bar{e}_{L}$ shared by all sentences of a language and sharpening semantic retrieval. Comparing centered alignment scores computed via Equation~\eqref{eq:mean-centering} against uncentered scores from Equation~\eqref{eq:mexa} quantifies how much of the measured alignment is carried by language-mean offsets rather than sentence-level semantics.

\begin{figure}[H]
\centering
\includegraphics[width=0.88\textwidth]{figures/fig_meancentering_geometry.pdf}
\caption{\textbf{Vector Space Geometry of Mean Centering}. 2D Principal Component Analysis (PCA) projection of high-dimensional sentence embeddings. (a) Uncentered sentence embeddings projected onto the first two principal components (PC1, PC2) are offset into distinct language clusters by a dominant language-identity vector $\bar{e}_{L}$. (b) Subtracting language means centers both manifolds at the origin, eliminating language identity bias and isolating semantic alignment.}
\label{fig:meancentering-geometry}
\end{figure}

\paragraph{Non-English pivot substitution.}
Equation~\eqref{eq:mexa} does not privilege English mathematically; the pivot $L_2$ is simply a parameter. We propose extending the pipeline to test the pivot hypothesis directly by recomputing all alignment scores across five diverse non-English pivot languages: high-resource Latin-script (French, German), non-Latin-script (Arabic, Chinese), and a language isolate (Basque). Both standard and centered variants are evaluated.

% ------------------------------------------------------------
\section{Models}
\label{sec:models}

We evaluate over 25 publicly available checkpoints, grouped into four categories that correspond to the four experimental axes. All models are obtained from the Hugging Face Hub and evaluated without fine-tuning; base (non-instruction-tuned) variants are used where available so that scores reflect pre-training rather than post-training.

\begin{table}[H]
\centering
\caption{\textbf{Evaluated language models}. Models are grouped by architectural category with parameter scale, vocabulary size, and training objective.}
\label{tab:evaluated-models-spec}
\small
\begin{tabular}{p{2.5cm}p{3.8cm}p{1.8cm}p{1.6cm}p{4.5cm}}
\toprule
\textbf{Category} & \textbf{Model Checkpoint} & \textbf{Parameters} & \textbf{Vocab Size} & \textbf{Training Objective / Focus} \\
\midrule
\textbf{Reproduction}
  & Llama 3.1 8B & 8.0B & 128,256 & Causal LM (Next-token pred.) \\
  & Mistral 7B v0.3 & 7.2B & 32,768 & Causal LM (Sliding window att.) \\
\midrule
\textbf{New Decoders}
  & Qwen 3 (0.6B / 1.7B / 4B / 8B) & 0.6B--8.0B & 151,646 & Multilingual Causal LM \\
  & Qwen 3.5 9B & 9.0B & 151,646 & Advanced Multilingual Causal LM \\
  & Apertus 8B / Apertus 1.5 8B & 8.0B & 32,000 & Open Swiss AI Multilingual LM \\
  & Apertus 1.5 70B & 70.0B & 32,000 & High-capacity Multilingual LM \\
\midrule
\textbf{Mixture-of-Experts}
  & Mixtral 8x7B / 8x22B & 46.7B / 141B & 32,000 & Sparse MoE (Top-2 routing) \\
  & Gemma 4 26B-A4B & 26.0B (4B act.) & 256,000 & Sparse MoE Causal LM \\
\midrule
\textbf{Encoders \& Embeddings}
  & XLM-RoBERTa (Base / Large) & 270M / 550M & 250,000 & Masked LM (Bidirectional) \\
  & Glot500 & 270M & 250,000 & Extended MLM (500+ languages) \\
  & LaBSE & 471M & 501,152 & Dual-Encoder Contrastive \\
  & mE5-base & 278M & 250,000 & Multilingual Text Embedding \\
  & Qwen 3 Embedding (0.6B--8B) & 0.6B--8.0B & 151,646 & Decoder-based Contrastive \\
\bottomrule
\end{tabular}
\end{table}

% ..............................................
\subsection{Reproduction Models}
\label{subsec:reproduction-models}

To validate our implementation against \citet{kargaran2025mexa}, we re-evaluate two models that have published reference scores: \textbf{Llama~3.1~8B} \citep{dubey2024llama3}, the primary English-centric decoder in the original paper, and \textbf{Mistral~7B~v0.3} \citep{jiang2023mistral}. If our $\mu_{\text{Max}}$/$\mu_{\text{Mean}}$ values agree with the published ones on the reference benchmark configurations (Section~\ref{sec:datasets}), the embedding extraction and scoring code can be considered correct, and the extended experiments can proceed on a verified foundation.

% ..............................................
\subsection{New Decoder Families}
\label{subsec:new-decoders}

The first extension applies MEXA to decoder families released after the original study.

Qwen~3 \citep{yang2025qwen3}, evaluated at 0.6B, 1.7B, 4B, and 8B parameters, was trained on a more multilingual data mixture than Llama. The size sweep isolates the effect of scale on alignment. Qwen~3.5~9B is the successor generation and lets us measure recipe-level progress at roughly constant parameter count. Apertus~8B and Apertus-Mini~4B are openly trained models with an explicitly multilingual, compliance-focused data policy, serving as a counterpoint to English-centric training. Gemma~4 spans five variants (E2B, E4B, 12B, 31B, and the sparse 26B-A4B), ranging from efficient edge-scale to 31B-parameter models.

\noindent
Together with the reproduction models, these give us size ladders within families (Qwen3 0.6B$\to$8B, Gemma~4 E2B$\to$31B) and recipe contrasts across families at matched scale (Llama~3.1~8B vs.\ Qwen3~8B vs.\ Apertus~8B).

% ..............................................
\subsection{Mixture-of-Experts Models}
\label{subsec:moe-models}

Sparse mixture-of-experts (MoE) architectures route each token through a small subset of expert feed-forward networks, so parameter count and per-token compute are decoupled \citep{jiang2024mixtral}. Whether expert routing helps or hurts cross-lingual alignment is an open question. Experts could specialize by language, which would fragment the shared space, or by function, which would leave alignment intact.

We evaluate \textbf{Mixtral~8x7B} and \textbf{Mixtral~8x22B} \citep{jiang2024mixtral} against the dense Mistral~7B~v0.3, and \textbf{Gemma~4~26B-A4B} against dense Gemma~4 variants. The MEXA pipeline applies unchanged in all cases, since it only consumes post-layer hidden states.

% ..............................................
\subsection{Encoder and Embedding Models}
\label{subsec:encoder-models}

The final model axis contrasts causal decoders with bidirectional encoders and purpose-built sentence-embedding models. The latter provide an upper bound for what alignment looks like when it is an explicit training objective rather than a by-product.

We evaluate three sub-categories. Masked-LM encoders include XLM-R base and large \citep{conneau2020xlmr}, Glot500 \citep{imani2023glot500} (an XLM-R variant extended to 500+ languages, particularly relevant for the Bible corpus tail), and mmBERT-base. Contrastively trained sentence encoders include LaBSE \citep{feng2022labse} and multilingual-E5-base \citep{wang2024multilingual}, both optimized so that translations are nearest neighbors in the embedding space. Decoder-based embedding models include the Qwen3-Embedding series at 0.6B, 4B, and 8B, which share the Qwen3 architecture but are fine-tuned for retrieval. This last group allows a controlled comparison of the same architecture before and after embedding-specific training.

% ------------------------------------------------------------
\section{Datasets}
\label{sec:datasets}

All experiments draw on the two parallel corpora introduced in Section~\ref{sec:parallel-corpora}. Configurations are chosen to match the original paper where reproduction is intended and to extend coverage otherwise. The two datasets differ in domain, vocabulary, and script distribution, which makes them complementary as a testbed.

\begin{table}[H]
\centering
\caption{\textbf{Parallel evaluation corpora and subsets}. Language coverage, sample size ($n$), domain, and role in the thesis.}
\label{tab:evaluation-datasets-overview}
\small
\begin{tabular}{p{3.5cm}lp{2.2cm}p{1.8cm}p{5.0cm}}
\toprule
\textbf{Corpus \& Subset} & \textbf{Languages} & \textbf{Sample Size ($n$)} & \textbf{Domain} & \textbf{Primary Thesis Purpose} \\
\midrule
\textbf{FLORES Benchmark} & 116 & 100 sentences & Wikipedia & Primary model ranking \& reproduction \\
\textbf{FLORES (1012)} & 116 & 1,012 sentences & Wikipedia & Distractor scaling ablation ($n$) \\
\textbf{FLORES Full} & 203 & 1,012 sentences & Wikipedia & Maximal high-quality coverage \\
\midrule
\textbf{Bible Benchmark} & 101 & 103 verses & Religious & Cross-domain stability check \\
\textbf{Bible Full} & >1,400 & Full text & Religious & Extremely low-resource tail \\
\bottomrule
\end{tabular}
\end{table}

\paragraph{FLORES-200.}
From the FLORES-200 devtest split \citep{nllbteam2022} we use three configurations. The FLORES benchmark subset contains the 116 languages that overlap with Belebele, with $n = 100$ parallel sentences per language. This is the exact configuration in \citet{kargaran2025mexa}, used for reproduction and all model-family comparisons. FLORES~(1012) uses the same 116 languages with all $n = 1{,}012$ devtest sentences, which lets us measure how sample size affects the retrieval criterion (larger $n$ means more distractors and a harder task). FLORES Full includes all 203 language-script pairs for maximal coverage.

FLORES-200 tags language-script pairs systematically (e.g., simplified Chinese \texttt{zho\_Hans}, traditional Chinese \texttt{zho\_Hant}, Cyrillic Russian \texttt{rus\_Cyrl}, Devanagari Hindi \texttt{hin\_Deva}), which lets us isolate script effects from language identity (for instance, Serbian in Cyrillic \texttt{srp\_Cyrl} vs.\ Latin \texttt{srp\_Latn}). The domain is contemporary Wikipedia prose covering geography, history, and science, with modern terminology and complex syntax.

\paragraph{Parallel Bible Corpus (sPBC).}
From the super-parallel Bible corpus \citep{mayer2014bible} we use two configurations. The Bible benchmark subset contains 101 languages with $n = 103$ parallel sentences, matching the original paper's setup. The Bible Full dataset covers the complete corpus of more than 1{,}400 language-script pairs, extending evaluation deep into the low-resource tail.

Unlike the modern encyclopedic prose of FLORES, the sPBC is entirely historical and religious, with archaic vocabulary, repetitive semantic motifs, and unusual grammatical constructs. It is also the most widely translated text available, spanning scripts absent from typical modern benchmarks: Ge'ez (Ethiopic) for Amharic and Tigrinya, Hebrew, Greek, and various minor regional scripts. For hundreds of low-resource languages that lack any digital crawl footprint or Wikipedia presence, the Bible corpus is the only viable means of evaluating representation alignment.

English (\texttt{eng\_Latn}) is the pivot in all standard runs. The pivot-substitution experiments of Section~\ref{sec:mexa-pipeline} use the FLORES benchmark subset.

% ------------------------------------------------------------
\section{Experimental Settings}
\label{sec:experimental-settings}

% ..............................................
\subsection{Evaluation and Analysis Methods}
\label{subsec:analysis-methods}

Unless stated otherwise, reported scores use position-weighted average embeddings (Equation~\eqref{eq:weighted-avg}) and are given as $\mu_{\text{Max}}$ and $\mu_{\text{Mean}}$ per model-dataset configuration; last-token embeddings and the centered variant are reported as ablations.

The analysis operates at four levels. At the model level, scores are averaged over all languages of a configuration for family and scale comparisons and for the reproduction check against published values. At the language level, per-language scores are ranked and grouped by script and resource level, then cross-referenced with tokenizer fertility (computed on both corpora) to separate tokenization effects from representational ones. At the layer level, full layer-wise alignment profiles $\mu C(\cdot, l)$ are examined to locate where in the network alignment peaks and to test the intermediate-layer prediction of the pivot hypothesis. At the validity level, Pearson correlations between MEXA scores and downstream accuracy on Belebele \citep{bandarkar2024belebele} and m-ARC \citep{lai2023okapi} are computed over the languages both benchmarks cover, testing whether alignment predicts task performance. Cross-model correlations measuring consistency across architectures are also evaluated.

For qualitative inspection, sentence embeddings are projected to two and three dimensions with PCA and t-SNE, colored by language and script, to visualize the clustering behavior that the retrieval score summarizes numerically.

% ..............................................
\subsection{Implementation and Compute}
\label{subsec:implementation}

The pipeline is implemented in Python with PyTorch and Hugging Face Transformers. The embedding extractor loads each checkpoint via \texttt{AutoModelForCausalLM} (or \texttt{AutoModel} for encoders) with output of all hidden states enabled, tokenizes each sentence with the model's own tokenizer, and stores per-layer sentence embeddings as serialized NumPy arrays, one file per language. The scoring module then builds the layer-wise cosine similarity matrices and evaluates Equation~\eqref{eq:mexa}.

Formatted per-language and per-layer scores are exported to CSV and JSON and explored in an interactive React dashboard developed for this thesis.

Small models ran locally; large-scale extraction jobs (8B parameters and above, MoE models, and full-corpus sweeps) were executed on the CIT-TUM-HN SLURM cluster at TUM Campus Heilbronn on NVIDIA GPUs.
